%=====================================================================
%  Fakerni — End-of-Studies Project Report (PFE)
%  Compile on Overleaf with pdfLaTeX.
%  Upload the whole "report/" folder (this file + the images/ subfolder).
%  Edit the \myXxx values just below with your real details.
%=====================================================================
\documentclass[12pt,a4paper,oneside]{report}

\usepackage[utf8]{inputenc}
\usepackage[T1]{fontenc}
\usepackage[english]{babel}
\usepackage[a4paper,margin=2.5cm]{geometry}
\usepackage{graphicx}
\usepackage{xcolor}
\usepackage{booktabs}
\usepackage{tabularx}
\usepackage{enumitem}
\usepackage{caption}
\usepackage{fancyhdr}
\usepackage{titlesec}
\usepackage{listings}
\usepackage[hidelinks]{hyperref}

\graphicspath{{images/}}

% ---- Your details (EDIT THESE) -------------------------------------
\newcommand{\myTitle}{Fakerni: A Smart Shopping and Task Coordination Platform for Groups}
\newcommand{\mySubtitle}{Full-Stack Web, Mobile and Real-Time Application with AI}
\newcommand{\myAuthor}{Ghassen Chich}          % <-- change to your full name
\newcommand{\mySupervisor}{[Supervisor Name]}   % <-- change
\newcommand{\myUniversity}{[University / Institute Name]}
\newcommand{\myDept}{[Department / Field of Study]}
\newcommand{\myYear}{2025 -- 2026}
% --------------------------------------------------------------------

% ---- Colors & code style -------------------------------------------
\definecolor{brand}{RGB}{0,54,61}
\definecolor{codebg}{RGB}{244,247,247}
\definecolor{codekw}{RGB}{26,107,115}
\definecolor{codecm}{RGB}{113,120,122}

\lstset{
  backgroundcolor=\color{codebg},
  basicstyle=\ttfamily\footnotesize,
  keywordstyle=\color{codekw}\bfseries,
  commentstyle=\color{codecm}\itshape,
  breaklines=true,
  frame=single,
  rulecolor=\color{lightgray},
  columns=fullflexible,
  showstringspaces=false,
  captionpos=b,
}

\titleformat{\chapter}[hang]
  {\normalfont\huge\bfseries\color{brand}}{\thechapter\quad}{0pt}{\Huge}
\titlespacing*{\chapter}{0pt}{0pt}{20pt}

\pagestyle{fancy}
\fancyhf{}
\fancyhead[L]{\small\color{brand}Fakerni}
\fancyhead[R]{\small\thepage}
\renewcommand{\headrulewidth}{0.4pt}

% Include an image only if it exists, otherwise print a placeholder box.
\newcommand{\safefig}[3]{% {file}{width}{caption}
  \begin{figure}[ht]\centering
  \IfFileExists{images/#1}{\includegraphics[width=#2]{#1}}%
  {\fbox{\parbox[c][3cm][c]{#2}{\centering\itshape Figure placeholder: upload \texttt{images/#1}}}}%
  \caption{#3}\label{fig:#1}
  \end{figure}}

\begin{document}

%=====================================================================
% TITLE PAGE
%=====================================================================
\begin{titlepage}
\centering
\vspace*{0.5cm}
{\large \myUniversity\par}
{\large \myDept\par}
\vspace{2cm}
{\Large\bfseries End-of-Studies Project Report\par}
\vspace{1.5cm}
\rule{\linewidth}{0.5mm}\\[0.5cm]
{\huge\bfseries\color{brand} \myTitle \par}
\vspace{0.4cm}
{\large\itshape \mySubtitle \par}
\rule{\linewidth}{0.5mm}\\[2cm]

\begin{tabular}{rl}
\textbf{Prepared by:} & \myAuthor\\[6pt]
\textbf{Supervised by:} & \mySupervisor\\[6pt]
\textbf{Academic year:} & \myYear\\
\end{tabular}

\vfill
{\small Source code: \url{https://github.com/ghassenchich/fakerni}\par}
\end{titlepage}

%=====================================================================
% ABSTRACT
%=====================================================================
\chapter*{Abstract}
\addcontentsline{toc}{chapter}{Abstract}
\thispagestyle{fancy}

Coordinating shared shopping and tasks within a group --- a family, a
community, an association or a society --- is commonly handled through informal
messaging and paper lists. This leads to duplicated purchases, forgotten items,
no accountability for who did what, and no visibility over shared spending.

This project presents \textbf{Fakerni}, a full-stack platform that gives any
group a shared, structured and \emph{real-time} space to organise what to buy
and do, available on both web and mobile. Beyond a conventional to-do
application, Fakerni introduces three differentiating capabilities:
\textbf{artificial intelligence} (turning a photographed receipt or a free-text
sentence into structured list items), \textbf{budget management and spending
analytics}, and \textbf{live collaboration} (instant synchronisation and
presence indicators). The system is built with a Django REST + Channels
backend, a React web client and a React Native (Expo) mobile client, is fully
trilingual (Arabic, French, English), and is validated by an automated test
suite of over 150 tests.

\vspace{0.4cm}
\noindent\textbf{Keywords:} group coordination, real-time web, Django, React,
React Native, WebSockets, artificial intelligence, budget analytics,
role-based access control.

%=====================================================================
% TOC
%=====================================================================
\tableofcontents
\listoffigures

%=====================================================================
\chapter{General Introduction}
%=====================================================================
Digital tools have transformed how individuals organise their daily lives, yet
the coordination of \emph{shared} responsibilities inside a group remains
surprisingly under-served. Families still coordinate grocery shopping over chat
messages; student clubs and associations track their errands on paper; nobody
has a reliable, shared view of what has been done or how much has been spent.

The aim of this end-of-studies project is to design and build a complete
platform, named \textbf{Fakerni} (from the Arabic \emph{fakkerni}, ``remind
me''), that addresses these shortcomings. A \emph{Fakra} (from the Arabic
\emph{fikra}, ``idea'') is a single shared list of items to buy or tasks to do;
a \emph{group} is the set of people who share it, with defined roles and an
invitation mechanism.

This report is organised as follows. Chapter~\ref{ch:context} presents the
context, the shortcomings of existing practices and the objectives of the
project. Chapter~\ref{ch:analysis} details the functional and non-functional
requirements. Chapter~\ref{ch:techarch} justifies the chosen technologies and
describes the overall architecture. Chapter~\ref{ch:design} covers the data
model, the authorisation model and the real-time design. Chapter~\ref{ch:impl}
describes the implementation of the main features. Chapter~\ref{ch:testing}
addresses testing and quality assurance, and Chapter~\ref{ch:deploy} discusses
deployment. A general conclusion then summarises the work and outlines
perspectives.

%=====================================================================
\chapter{Project Context and Problem Statement}\label{ch:context}
%=====================================================================
\section{Context}
The project sits at the intersection of two mature domains: collaborative task
management (to-do and list applications) and shared expense management. Popular
task applications excel at individual productivity but rarely model a
\emph{group} as a first-class entity with roles and real-time collaboration.
Conversely, expense-splitting applications track shared money but do not manage
the underlying shopping or task lists. Fakerni unifies both around a single
shared object, the Fakra.

\section{Study of the Existing Situation}
Current group coordination typically relies on:
\begin{itemize}[leftmargin=1.4em]
  \item \textbf{Instant messaging} (e.g.\ group chats): unstructured, items get
        lost in the conversation, and there is no notion of ``done''.
  \item \textbf{Paper or personal note apps}: not shared; only one person sees
        the list.
  \item \textbf{Generic to-do apps}: individual-centric, without group roles,
        shared budgets or live synchronisation.
\end{itemize}

\section{Problem Statement}
From the study above, four concrete problems emerge:
\begin{enumerate}[leftmargin=1.6em]
  \item \textbf{Scattered information} --- no single shared source of truth.
  \item \textbf{Duplication and omissions} --- two people buy the same thing,
        or an unclaimed item is forgotten.
  \item \textbf{No accountability} --- no record of who added or completed what.
  \item \textbf{No spending awareness} --- no shared view of cost or budget.
\end{enumerate}

\section{Proposed Solution and Objectives}
Fakerni proposes a shared, real-time platform providing: a single source of
truth (groups, Fakras, items); instant synchronisation across web and mobile;
an activity log for accountability; per-item pricing with budgets and analytics;
and AI assistance to reduce manual data entry. The measurable objectives were:
a secure multi-user REST + real-time API, two client applications (web and
mobile) sharing the same backend, trilingual support, and a comprehensive
automated test suite.

%=====================================================================
\chapter{Requirements Analysis}\label{ch:analysis}
%=====================================================================
\section{Actors}
\begin{itemize}[leftmargin=1.4em]
  \item \textbf{Visitor} --- can register and log in.
  \item \textbf{Member} --- a registered user who belongs to a group; can create
        and manage Fakras and items.
  \item \textbf{Admin / Owner} --- a member with elevated privileges over a group
        (manage members, moderate content, delete lists).
\end{itemize}

\section{Functional Requirements}
The main functional requirements are summarised in Table~\ref{tab:func}.

\begin{table}[ht]\centering
\caption{Main functional requirements}\label{tab:func}
\begin{tabularx}{\linewidth}{@{}lX@{}}
\toprule
\textbf{Module} & \textbf{Capabilities}\\
\midrule
Accounts & Registration, JWT login, password reset by e-mail OTP, profile, biometric app-lock (mobile).\\
Groups & Types (family/community/organization/society), invitation codes, roles (owner/admin/member), member management.\\
Fakras & Personal, shared and recurring lists; archive, duplicate, search and filter.\\
Items & Create/edit, mark done/undo (10-minute window), quantity, unit, category, price, photo attachments, assignment to a member.\\
AI & Smart Add (text $\rightarrow$ items), Smart Scan (receipt $\rightarrow$ items+prices), Suggestions, natural-language Commands, spending digest.\\
Budgets & Per-Fakra budget, progress and over-budget alerts, spending analytics by month/category/member, expense splitting (``settle up'').\\
Real-time & Live synchronisation and presence (who is viewing a list).\\
Notifications & Push notifications for new items, completions, members, reminders and budget alerts.\\
Exports & CSV (both platforms) and server-generated PDF.\\
\bottomrule
\end{tabularx}
\end{table}

\section{Non-Functional Requirements}
\begin{itemize}[leftmargin=1.4em]
  \item \textbf{Security:} hashed passwords, short-lived JWT, rate limiting,
        encryption of sensitive data at rest, centralised authorisation.
  \item \textbf{Performance:} optimised database queries, stateless API enabling
        horizontal scaling.
  \item \textbf{Usability:} responsive interfaces, dark mode, trilingual UI with
        right-to-left support for Arabic.
  \item \textbf{Reliability:} automated tests and continuous integration.
  \item \textbf{Portability:} containerised deployment (Docker).
\end{itemize}

%=====================================================================
\chapter{Technologies and Architecture}\label{ch:techarch}
%=====================================================================
\section{Technology Choices}
Table~\ref{tab:tech} lists the main technologies and the rationale behind each
choice.

\begin{table}[ht]\centering
\caption{Technology stack and justification}\label{tab:tech}
\begin{tabularx}{\linewidth}{@{}lX@{}}
\toprule
\textbf{Layer} & \textbf{Technology and justification}\\
\midrule
Backend framework & \textbf{Django + Django REST Framework} --- mature ORM, authentication, serializers and throttling, enabling a correct and secure API to be built quickly.\\
Real-time & \textbf{Django Channels + Daphne (ASGI)} --- adds WebSockets within the same codebase and authentication system, avoiding a separate service.\\
Database & \textbf{PostgreSQL} --- relational integrity for the users/groups/lists/items hierarchy, transactions and good concurrency.\\
Authentication & \textbf{JWT (SimpleJWT)} --- stateless tokens allowing the API to scale horizontally without server-side sessions.\\
Artificial intelligence & \textbf{Google Gemini} --- fast, low-cost model called with structured output schemas to guarantee valid JSON responses.\\
Web client & \textbf{React 19 + Vite + Tailwind CSS} --- fast development and a consistent component-based UI.\\
Mobile client & \textbf{React Native + Expo} --- a single JavaScript codebase for Android and iOS, reusing the web application's logic.\\
Infrastructure & \textbf{Docker, docker-compose, GitHub Actions} --- reproducible environment and continuous integration.\\
\bottomrule
\end{tabularx}
\end{table}

\section{Global Architecture}
Fakerni follows a classic three-tier architecture with an added real-time
channel, illustrated in Figure~\ref{fig:architecture.png}. The two clients (web
and mobile) communicate with the backend through a REST API over HTTPS and a
WebSocket connection over WSS. The backend, served by the Daphne ASGI server,
exposes the REST endpoints and the WebSocket consumers, and connects to
PostgreSQL for persistence, to Redis for the real-time channel layer, and to
external services (Gemini for AI, Firebase Cloud Messaging for push
notifications).

\safefig{architecture.png}{\linewidth}{Three-tier, real-time system architecture.}

Because the REST API is stateless (authentication is carried by a JWT on each
request), it can be replicated horizontally behind a load balancer; the Redis
channel layer then fans real-time events out across all instances.

%=====================================================================
\chapter{Design}\label{ch:design}
%=====================================================================
\section{Data Model}
The core entities and their relationships are shown in the
entity--relationship diagram in Figure~\ref{fig:erd.png}. A \texttt{User}
belongs to \texttt{Household}s (the generic ``group'' entity) through
\texttt{Membership}, which carries the role. A \texttt{Fakra} may belong to a
group or be personal (its group reference is nullable). A \texttt{Fakra}
contains \texttt{Item}s, which may have \texttt{ItemAttachment}s and an
assignee; every action is recorded in the \texttt{ActivityLog}; and a personal
Fakra can be explicitly shared through \texttt{FakraAccess}.

\safefig{erd.png}{\linewidth}{Entity--relationship diagram of the core data model.}

\section{Authorisation Model}
A distinctive design decision is that access to a Fakra can be granted in three
ways: the user created it, the user belongs to its group, or it was explicitly
shared with them. All authorisation decisions --- who may view, edit, delete,
archive or share a Fakra, and who may modify or assign an item --- are
centralised in a single module (\texttt{fakras/permissions.py}). Views never
re-implement rules inline; they call helper functions. The API additionally
returns, for each Fakra, the set of actions the current caller is permitted to
perform, so the clients can hide forbidden actions from the interface.

\section{Real-Time Design}
The real-time layer uses three channel scopes: a per-group channel, a per-Fakra
channel (which also works for personal and shared lists), and a per-user
channel used to keep each user's dashboard current. Because JSON Web Tokens are
not automatically applied to WebSocket connections, a custom ASGI middleware
validates the token from the connection query string and attaches the
authenticated user to the connection; each consumer then verifies authorisation
before joining a channel group, which prevents a user from listening on a
channel that is not theirs.

%=====================================================================
\chapter{Implementation}\label{ch:impl}
%=====================================================================
This chapter describes the realisation of the main features. The backend is
organised into three Django applications: \texttt{users} (accounts and
identity), \texttt{household} (groups and roles) and \texttt{fakras} (lists,
items, AI, budgets and real-time).

\section{Authentication and Accounts}
Authentication is token-based. On login the server verifies the password
against a salted PBKDF2 hash and issues a short-lived access token and a
longer-lived refresh token. The clients attach the access token to every
request and transparently refresh it when it expires. Figure~\ref{fig:web-login.png}
shows the web login screen.

\safefig{web-login.png}{0.7\linewidth}{Web authentication screen.}

\section{Groups, Fakras and Items}
Users create groups and invite others through an expiring invitation code.
Within a group, any member may create a shared Fakra; editing or deleting a
Fakra, or moderating items created by others, is governed by the centralised
authorisation model. Items support quantities, units, categories, prices, photo
attachments, a ten-minute ``undo'' window after completion, and assignment to a
specific member. Figure~\ref{fig:web-dashboard.png} shows the web dashboard and
Figure~\ref{fig:web-fakra.png} a Fakra detail view.

\safefig{web-dashboard.png}{\linewidth}{Web dashboard with the overview strip and list of Fakras.}
\safefig{web-fakra.png}{\linewidth}{Fakra detail: items, budget, AI Smart Add and export actions.}

\section{Artificial Intelligence}
Four AI capabilities are built on Google Gemini, each called with a strict
output schema so that the model returns valid structured data rather than free
text: \emph{Smart Add} converts a sentence such as ``milk, eggs and 2kg rice''
into structured items; \emph{Smart Scan} reads a receipt photograph and extracts
items together with their prices; \emph{Suggestions} proposes relevant
additional items; and \emph{Smart Command} interprets natural-language commands
(``mark milk as done, remove the bread''). A fifth feature, the \emph{spending
digest}, produces a short natural-language summary of a group's spending. If the
AI service is unavailable, the corresponding endpoints fail gracefully and the
rest of the application continues to work.

\section{Budgets, Analytics and Expense Splitting}
Each item can carry an estimated price. A Fakra may define an explicit budget,
against which the application computes the amount spent, the remaining amount and
a progress indicator, and raises an alert when the budget is exceeded. A
dedicated analytics endpoint aggregates spending by month, by category and by
member, and supports time-range filtering. Finally, an expense-splitting feature
computes, from the ``who bought what'' data, each member's balance relative to
their fair share and proposes a minimal set of transfers to settle the group's
accounts.

\section{Real-Time Collaboration}
Changes to a list --- adding an item, marking it done, uploading an attachment
--- are broadcast instantly to all connected members. A presence feature shows,
in real time, which members are currently viewing a given Fakra. On the mobile
client, the same backend is complemented by native capabilities: camera-based
barcode scanning, biometric application lock and native push notifications.
Figures~\ref{fig:mobile-dashboard.png} and~\ref{fig:mobile-fakra.png} show the
mobile client.

\safefig{mobile-dashboard.png}{0.32\linewidth}{Mobile dashboard.}
\safefig{mobile-fakra.png}{0.32\linewidth}{Mobile Fakra detail.}

\section{Internationalisation}
The two clients are fully translated into English, French and Arabic. Selecting
Arabic switches the whole interface to a right-to-left layout, not merely the
text.

%=====================================================================
\chapter{Testing and Quality Assurance}\label{ch:testing}
%=====================================================================
Quality was addressed continuously rather than as a final step. The backend is
covered by an automated suite of more than \textbf{150 tests} spanning the three
applications, including authentication flows, the centralised authorisation
rules, the ten-minute undo window, budget alerts, the AI endpoints (with the
external model mocked), the expense-splitting computations, and the WebSocket
delivery of real-time events (including presence on join and leave). The web
client additionally has unit tests for its utility functions.

A continuous-integration pipeline (GitHub Actions) runs the full backend suite
against a PostgreSQL service, together with the web unit tests and production
build, on every push. The OpenAPI schema of the API is generated automatically
and is available through interactive Swagger and Redoc interfaces.

%=====================================================================
\chapter{Deployment}\label{ch:deploy}
%=====================================================================
The application is fully containerised. A \texttt{Dockerfile} builds the backend
image, and a \texttt{docker-compose} configuration orchestrates the database,
the Redis channel layer and the application server. A deployment blueprint for
the Render platform is also provided, allowing the backend and its managed
PostgreSQL database to be provisioned automatically; the web client is built as
a static bundle and served by any static host. Configuration is entirely driven
by environment variables, and production security settings (HTTPS redirection,
secure cookies, HSTS) are enabled automatically when debugging is disabled.

%=====================================================================
\chapter*{Conclusion and Perspectives}
\addcontentsline{toc}{chapter}{Conclusion and Perspectives}
%=====================================================================
This project resulted in a complete, tested and deployable platform that turns
the informal coordination of a group's shopping and tasks into a shared,
real-time and financially aware experience, available on both web and mobile.
The main contributions are a unified data model centred on the Fakra, a
centralised and spoof-proof authorisation model, an integrated artificial-
intelligence layer that removes manual data entry, a complete budgeting and
analytics module including expense splitting, and a multi-scope real-time layer
with presence.

Several perspectives would extend the work further. A history-aware
\emph{predictive restock} engine and \emph{price-anomaly} detection are already
partially in place and could be deepened. Collaboration could be enriched with
comments and typing indicators; the AI could generate a shopping list directly
from a meal plan; and the mobile client could gain an offline mode with
background synchronisation and location-based reminders. These directions build
naturally on the foundations established in this project.

\end{document}
